%% notes-tex/025-additive-baselines/sections/5-ladder.tex

\section{The ladder under both splits\secstatus{todo}}
\label{sec:ladder}

%% SOURCE: experiments/025-solid-growth/scripts/additive_baselines_025_panels.py ->
%% tables/t2-heldout.tex, from results/additive_baselines_025.csv, the summary json,
%% and 010-kuzmin-tmi/results/additive_baseline_gene_interaction.csv.
\begin{table}[htbp]
\centering
\footnotesize
\caption{Test Pearson for every model under both arms, beside the 010 numbers. B0
through B4 are closed-form fits on a fixed split and carry no spread; B5 is the mean
and standard deviation over three seeds. The transformer rows stand on different
footings, named in the row label: the 010 checkpoints were scored on the 010 test
split of 37{,}673 records; the GH 1598 and GH 1640 validation maxima are over 36
logged epochs; the last row is GH 1640's validation-selected checkpoint scored on
the arm Q test part.}
\label{tab:heldout}
%% GENERATED by experiments/025-solid-growth/scripts/additive_baselines_025_panels.py
%% SOURCE: the result files named in that script's docstring. Do not edit by hand.
\begin{tabular}{lrrr}
\toprule
 & \multicolumn{2}{c}{Arm R, random over records} & Arm Q, query-pair disjoint \\
Model & 010 build & 025 build & 025 build \\
\midrule
B0 train mean & 0.000 & 0.000 & 0.000 \\
B4 query pair only & 0.390 & 0.390 & 0.000 \\
B3 screen mean & 0.390 & 0.390 & 0.142 \\
B1 additive per-gene ridge & 0.400 & 0.400 & 0.185 \\
B2 additive plus pair ridge & 0.406 & 0.406 & 0.180 \\
B5 embedding MLP, 3 seeds & $0.426 \pm 0.001$ & $0.432 \pm 0.005$ & $0.141 \pm 0.007$ \\
CGT, three 010 checkpoints, test & 0.438 to 0.455 & & \\
CGT GH 1598, val max, epoch 14 of 36 & & 0.446 & \\
CGT GH 1640, val max, epoch 7 of 36 & & & 0.199 \\
\bottomrule
\end{tabular}

\end{table}

%% SOURCE: experiments/025-solid-growth/scripts/additive_baselines_025_panels.py ->
%% additive_baselines_025_fig1_ladders.svg. Panels a, b: the test rows of
%% results/additive_baselines_025.csv plus the transformer references in the summary
%% json. Panel c: results/additive_baselines_025_arm_val_history.csv. Panel d:
%% 010-kuzmin-tmi/results/query_pair_disjoint_cv.csv against the arm Q test rows.
\tcfig{figures/additive_baselines_025_fig1_ladders.pdf}{The ladder under both
splits. \textbf{a}, arm R, random over records: the six baselines on the 37{,}673
test records, the three 010 checkpoints on the 010 test split, and the 025
replication GH 1598 as a validation maximum. \textbf{b}, arm Q, query-pair
disjoint, on the 37{,}791 test records of 46 held-out screens, with GH 1640 as a
validation maximum and its epoch 7 checkpoint scored on test. Within each arm, a
light bar is a test score and a dark bar a validation maximum. Only B5 carries an
error bar, the standard
deviation over its three seeds; every other bar is one deterministic fit.
\textbf{c}, validation Pearson per epoch for the two arms of the 010 configuration,
each against the additive ridge fit on its own validation part; open circles mark
each run's best epoch, the value the dark bars of a and b carry. \textbf{d}, how far
the disjoint null moves with the choice of held-out screens: five disjoint folds on
the 010 build per model, mean and standard deviation with each fold as a point, and
the arm Q single split as a diamond. B0 and B4 are zero on every fold and are
omitted.}{fig:ladders}

\subsection{Reading the ladder\secstatus{todo}}

\notesec{Arm R is the 010 result}
Every closed-form model agrees with its 010 value to $3 \times 10^{-13}$, which
follows from the record-level label and split identity of Section~\ref{sec:data}
rather than testing it. B5 differs by 0.006 against a three-seed standard deviation
of 0.005; the 025 fit ran on CPU where 010's ran on GPU, and the cause is untraced.
GH 1598 reaches a validation maximum of 0.446 at epoch 14, within 0.001 of the lower
end of the 0.447 to 0.462 band the three 010 checkpoints recorded on validation.

\notesec{The additive null is close, but the transformer does beat it}
B1 reaches 0.400 against the transformer's 0.438 to 0.455, so a model that cannot
represent any interaction recovers 88 to 91 percent of the correlation. In explained
variance the gap is $r^2 = 0.160$ against $0.207$ for the best checkpoint, an
increment of 0.047 of label variance. This is not the DeWitt result, where the
network collapsed onto the additive model at $r > 0.999$. The margin is real,
$+0.038$ to $+0.055$ with paired-bootstrap intervals that exclude
zero\iflong{} (Section~\ref{sec:paired})\fi, and it is small.

\iflong
\notesec{Seed spread is a large fraction of that margin}
The three transformer runs span 0.438 to 0.455, a standard deviation of 0.0088.
The three B5 seeds span 0.425 to 0.427 on the 010 build, a standard deviation of
0.0011. So the transformer's run-to-run spread is eight times the nonlinear
baseline's and roughly a fifth of its own margin over B1. A single checkpoint's
number is not the capability, and the three runs are not seed replicates in the
strict sense: the training script fixes the seed at 42, M01 and M02 differ only in
run-to-run nondeterminism, and M03 additionally changes the scheduler's first
cycle length.

\notesec{Capacity closes about half the gap, and the graphs cannot be ruled out}
B5 sees exactly B1's features, has no graph and no attention, and reaches 0.426 on
the 010 build and 0.432 on 025. That closes about half the distance from B1 to the
transformer. Attributing the remainder to capacity alone does not hold, because the
graphs shape the trained weights even though they are absent from the prediction
function (Section~\ref{sec:graphs}), and Section~\ref{sec:graphablation} measures
how heavily. What the ladder alone says is that richer set aggregation and eight
times the parameters are worth at most 0.03 Pearson over B5, and that how much of
that 0.03 is the graphs is not separable from these runs.

\notesec{Spearman compresses the ladder further}
On the 010 build the rank correlations run 0.406 for B1, 0.412 for B2 and 0.417 to
0.421 for B5, against 0.424 to 0.426 for the transformer. The gap from the additive
null to the transformer is 0.020 in Spearman against 0.045 in Pearson. Ranking is
what a design-space selection consumes, so the smaller of the two is the more
decision-relevant number.
\fi

\notesec{On the random split the ladder is mostly screen identity}
B4 knows only which query pair a triple carries and never looks at the third gene.
It reaches 0.390, 97 percent of the additive ridge and 86 percent of the best
checkpoint. Those 420 pairs are the Kuzmin query doubles, the most frequent
appearing 3{,}308 times, so B4 is close to a per-screen offset, and on a split that
is random over records most of what every model in Table~\ref{tab:heldout} is
scoring is the identity of the screen a triple came from. It also explains why the
additive null does so well on a label constructed to have the additive part
removed: what B1's per-gene coefficients absorb is largely query-pair structure
reaching them through the genes that make up the query pairs.

\notesec{Holding out screens collapses the ladder and inverts it}
On arm Q the additive ridge falls from 0.400 to 0.185 and B4 to zero by
construction. B5 falls from 0.432 to 0.141, below the ridge: the capacity that
helped on unseen records of seen screens hurts on unseen screens. B3 falls back from
a screen mean to a gene mean and lands at 0.142. The one transformer checkpoint
with a test score on this split, GH 1640 at epoch 7, lands at 0.139, level with B5
and below the ridge (Section~\ref{sec:disjoint}).

\notesec{One disjoint split is a small sample of screens}
Panel d is the reason arm Q's numbers are read as one draw rather than as the
disjoint null. Across five disjoint folds on the 010 build the ridge spans 0.088 to
0.151, $0.127 \pm 0.033$, and arm Q's 0.185 lies above every fold. The same holds
for B2, B3 and B5, so arm Q is a comparatively easy draw of 46 held-out screens. A
transformer number on arm Q is compared to the arm Q nulls, on the same 46 screens,
and not to the fold mean.

\iflong
%% SOURCE: experiments/010-kuzmin-tmi/scripts/query_pair_disjoint_split.py and
%% query_pair_disjoint_cv.py -> results/query_pair_disjoint_cv.csv and
%% additive_baseline_gene_interaction_query_pair_disjoint.csv
\begin{table}[htbp]
\centering
\footnotesize
\caption{Held-out Pearson on the 010 build under the published random-over-records
split and under splits that never share a Kuzmin query double. The single disjoint
split is 010's own (285, 67 and 68 query pairs), a different partition of the same
420 pairs from arm Q's; the cross-validated column is the mean over five folds, each
holding out a disjoint set of query pairs, with B5 additionally averaged over three
seeds. B4 is structurally zero once query pairs are disjoint.}
\label{tab:disjoint}
\begin{tabular}{lrrr}
\toprule
& \multicolumn{1}{c}{Random split} & \multicolumn{2}{c}{Query-pair disjoint} \\
Model & test $r$ & one split & 5-fold mean $\pm$ sd \\
\midrule
B0 train mean & 0.000 & 0.000 & $0.000 \pm 0.000$ \\
B4 query pair only & 0.390 & 0.000 & $0.000 \pm 0.000$ \\
B3 hierarchical mean & 0.390 & 0.173 & $0.066 \pm 0.039$ \\
B1 additive per-gene ridge & 0.400 & 0.174 & $0.127 \pm 0.033$ \\
B2 additive plus pair ridge & 0.406 & 0.170 & $0.124 \pm 0.034$ \\
B5 embedding MLP & 0.426 & 0.150 & $0.058 \pm 0.037$ \\
\bottomrule
\end{tabular}
\end{table}

Under the disjoint folds the nonlinear baseline early-stops at epoch 3 to 7 rather
than epoch 20 to 30, having reached training Pearson 0.43 while held-out
performance was already falling. Generalizing to an unseen query double is a
different problem from generalizing to an unseen record of a seen query double, and
the extra capacity that helped on the second hurts on the first.

\subsection{Verifying the architecture reading\secstatus{todo}}
\label{sec:verify}

Section~\ref{sec:sees} claims the encoder is strain-independent and that no
adjacency tensor is an input to the prediction. Both are checkable rather than
merely readable, and checking them also produces the per-record predictions that
Section~\ref{sec:paired} needs. If the encoder really does not see the strain, then
scoring a record needs only the perturbed gene indices, and the index space is the
sorted base-graph node list, rebuildable from the genome without touching the
LMDB. \file{experiments/010-kuzmin-tmi/scripts/score_010_checkpoints_directly.py}
does that, loads each 010 checkpoint into the model class, and recomputes the
metrics. The check is that the recomputed numbers must reproduce what the original
evaluation runs recorded: a wrong gene ordering or a mis-loaded weight would give a
plausible but wrong number.

%% SOURCE: experiments/010-kuzmin-tmi/scripts/score_010_checkpoints_directly.py
%% -> results/cgt_direct_scoring.csv, against the recorded eval-run summaries
\begin{table}[htbp]
\centering
\footnotesize
\caption{Scoring the 010 checkpoints through a hand-built forward pass that
computes the encoder once and feeds it only perturbed-gene indices reproduces the
recorded metrics. The residual difference is at the fourth decimal and is
consistent with the original runs using mixed bf16 precision.}
\label{tab:repro}
\begin{tabular}{lrrrr}
\toprule
& \multicolumn{2}{c}{Val Pearson} & \multicolumn{2}{c}{Val MSE} \\
Checkpoint & recomputed & recorded & recomputed & recorded \\
\midrule
M01 \texttt{lzs9pcj3} & 0.451970 & 0.451963 & 0.003244 & 0.003244 \\
M02 \texttt{yv4r30bi} & 0.447319 & 0.447155 & 0.003258 & 0.003259 \\
M03 \texttt{c7671wgj} & 0.461917 & 0.461881 & 0.003163 & 0.003163 \\
\bottomrule
\end{tabular}
\end{table}

Two details had to be right, and getting either wrong produced a plausible but
wrong answer. The index space is the sorted genome gene set, 6{,}607 genes, not the
build's own 6{,}579 gene set; using the latter gave a validation Pearson of $-0.001$
to $0.002$. And the readout pools by mean, whereas the current class defaults to
sum; using the default gave 0.420 to 0.438, close enough to look credible and wrong.
That reproduction is the direct verification of Section~\ref{sec:encoder}: the
encoder was evaluated exactly once, on a strain-independent input, and the recorded
per-split metrics came back. Setting the graph penalty weight from $1.0$ to $0.0$
on a trained checkpoint changes the maximum absolute test prediction by
$9.0 \times 10^{-7}$, $1.2 \times 10^{-6}$ and $1.3 \times 10^{-6}$, against a
prediction spread of about $0.036$: adjacency is not an input to the prediction
function. It says nothing about whether the graphs shaped the weights being scored.

\subsection{What the graphs did during training\secstatus{todo}}
\label{sec:graphablation}

Three independent measurements, none of which the inference-time ablation can see.

\notesec{The penalty was almost the whole objective}
The trainer logs the graph term, the point term and their normalized shares.
Against a validation point loss of $0.81$ the graph term is $5{,}740$, and the
share \texttt{norm\_weighted\_graph\_reg} is $0.99986$ for all three checkpoints
(Figure~\ref{fig:penalty}). Eq.~\ref{eq:effcoef} accounts for the size.

\notesec{The regularized heads recover their graphs}
The trainer scores each regularized head by asking, for every gene $i$ with degree
$d_i$, what fraction of $i$'s true neighbors fall in that head's top $d_i$ attention
targets.

%% SOURCE: val_edge_recovery/{graph}_L1_H{k}/recall_at_deg in the three
%% checkpoints' re-evaluation runs under $DATA_ROOT/wandb-experiments.
\begin{table}[htbp]
\centering
\footnotesize
\caption{Neighbor recall at each gene's own degree, layer 1, one head per graph.
Ordinary self-attention has no reason to place its top $d_i$ mass on a gene's true
neighbors, so these values are attributable to the penalty of
Eq.~\ref{eq:graphreg}.}
\label{tab:edgerecovery}
\begin{tabular}{llrrr}
\toprule
Graph & Head & M01 & M02 & M03 \\
\midrule
physical & 0 & 0.940 & 0.940 & 0.938 \\
regulatory & 1 & 0.992 & 0.992 & 0.991 \\
tflink & 2 & 0.979 & 0.979 & 0.979 \\
string neighborhood & 3 & 0.952 & 0.951 & 0.950 \\
string fusion & 4 & 0.999 & 1.000 & 1.000 \\
string cooccurence & 5 & 1.000 & 1.000 & 1.000 \\
string coexpression & 6 & 0.729 & 0.731 & 0.729 \\
string experimental & 7 & 0.774 & 0.772 & 0.773 \\
string database & 8 & 0.996 & 0.996 & 0.995 \\
\bottomrule
\end{tabular}
\end{table}

\notesec{Removing the penalty stops the model training}
A matched pair of 30-epoch configurations exists for exactly this test, identical
but for the graph term's weight. With the weight at $1.0$, three replicates reach
validation Pearson $0.464$, $0.456$ and $0.456$. With it at $0.0$, the three
companion runs reach $-0.007$, $-0.031$ and a diverged run that returned no
number, with training Pearson around $0.002$. That establishes that the graph term
is load-bearing for optimization in this configuration. It does not establish that
the graphs supply biological information the model uses, because a term carrying
99.99 percent of the loss is also doing whatever conditioning, scaling and implicit
regularization that share implies. Separating the two would need a control with a
comparable penalty against permuted or degree-matched random graphs, which has not
been run.

\subsection{Record-by-record agreement\secstatus{todo}}
\label{sec:paired}

With per-record predictions available, the comparison can be made the way the
DeWitt preprint makes it, between the predictions themselves rather than between
their summary statistics. These are 010 checkpoints on the 010 test split, which is
arm R's test part record for record.

%% SOURCE: experiments/010-kuzmin-tmi/scripts/paired_prediction_agreement.py
%% -> results/paired_prediction_agreement.csv
\begin{table}[htbp]
\centering
\footnotesize
\caption{Correlation between models' test-set predictions, 37{,}673 records. The
transformer does not reproduce the additive model, which is where 010 departs from
the MULTI-evolve case. But two transformer training runs agree with each other at
0.788 to 0.821, barely more than a transformer agrees with the additive model at
0.726 to 0.759, and no more than it agrees with the nonlinear baseline at 0.784 to
0.810.}
\label{tab:paired}
\begin{tabular}{lrrr}
\toprule
& \multicolumn{3}{c}{Pearson with} \\
Model & CGT M01 & CGT M02 & CGT M03 \\
\midrule
B1 additive ridge & 0.752 & 0.726 & 0.759 \\
B2 additive plus pair & 0.759 & 0.734 & 0.765 \\
B4 query pair only & 0.729 & 0.706 & 0.735 \\
B5 embedding MLP & 0.788 & 0.784 & 0.810 \\
CGT M01 & & 0.788 & 0.821 \\
CGT M02 & 0.788 & & 0.804 \\
\bottomrule
\end{tabular}
\end{table}

\tcfig{figures/paired_prediction_agreement.pdf}{Test-set predictions of one model
against another, 37{,}673 records as a log-density hexbin. \textbf{a}, the best
transformer checkpoint against the additive ridge. \textbf{b}, against the
nonlinear baseline on the same features. \textbf{c}, against a second transformer
training run. The transformer is no closer to a second training run of itself than
it is to a nonlinear baseline that has no graph and an eighth of the
parameters.}{fig:paired}

\notesec{The headline distinction from the preprint}
DeWitt's networks correlated with the additive model at $r > 0.999$, so the network
\emph{was} the additive model. Here the correlation is $0.73$ to $0.76$, so the
transformer is computing something the additive model does not.

\notesec{How reproducible the extra content is}
Regressing each model's predictions on B1's and correlating what is left isolates
the non-additive component. It is 43.5, 47.3 and 42.4 percent of each checkpoint's
prediction variance. Between two checkpoints it correlates at $0.533$, $0.584$ and
$0.564$. Against the part of the label B1 also fails to explain it correlates at
$0.236$, $0.234$ and $0.253$, where B5's non-additive component manages $0.165$. So
the non-additive content is neither noise nor stable: about half of it is shared
between two training runs, and the shared part carries real signal about the
residual label.

\notesec{The models are wrong in the same places}
Correlating the residuals $y - \hat{y}$ of B1 and each checkpoint gives $0.907$,
$0.897$ and $0.918$. Almost all of the error is shared, so the transformer is not
fixing a class of records the additive model gets wrong; it is making a modest,
broadly distributed improvement.

\notesec{The advantage is real and bounded}
A paired bootstrap over test records, 2{,}000 resamples, puts the Pearson gain over
B1 at $+0.043$ with a 95 percent interval $[+0.023, +0.063]$ for M01, $+0.038$
with $[+0.019, +0.058]$ for M02, and $+0.055$ with $[+0.034, +0.074]$ for M03. All
three intervals exclude zero. B5 gains $+0.025$ with $[+0.017, +0.033]$.

\notesec{Rankings diverge where selection happens}
Among the 100 records each model calls most negative, B1 and M03 share 31, B5 and
M03 share 45, and B1 and B5 share 46. At $K = 10{,}000$ all pairs exceed 0.91.
Between two transformer checkpoints the overlap rises monotonically with $K$: 0.10
to 0.30 at $K = 10$, 0.39 to 0.47 at $K = 100$, 0.57 to 0.60 at $K = 1{,}000$, and
0.91 at $K = 10{,}000$. An extreme quantile is where the fitted surface is least
constrained by data and where the label itself is noisiest, so two runs that agree
about the bulk should be expected to disagree about the last hundred; the
consequence for practice is a reason to ensemble.

\notesec{Prediction spread}
Test prediction standard deviations are 0.025 for B1, 0.029 for B5 and 0.036 to
0.038 for the checkpoints, against a test label standard deviation of 0.064. Every
model is heavily shrunk toward the mean, as a squared-error fit on a noisy target
should be.
\fi
